% !TEX root = ../../thesis.tex

%%%%%%%%%%%%%%%%%%%%%%%%%%%%%%%%%%%%%%%%%%%%%%%%%%%%%%%%%%%%%%%%%%%%%%%%%%%%%%%
\section{ONVP Sweepouts}\label{sec:non-excessive-prelim}
%%%%%%%%%%%%%%%%%%%%%%%%%%%%%%%%%%%%%%%%%%%%%%%%%%%%%%%%%%%%%%%%%%%%%%%%%%%%%%%

We now review the non-excessive min-max framework of Chodosh--Liokumovich--Spolaor~\cite{CLS22}. The central idea is to refine the class of optimal sweepouts by imposing a condition---\emph{non-excessiveness}---that prevents critical slices from being locally replaceable by competitors of strictly smaller mass. This additional structure, absent in the classical Almgren--Pitts theory, leads to strong geometric control on the limit varifold. We work throughout with ONVP sweepouts (Definition~\ref{def:p2-onvp}).

Given an optimal sweepout, the \emph{critical domain} identifies the parameter values at which the mass is maximal. The non-excessive condition asks that these critical parameters cannot be avoided by local replacements.

\begin{definition}[Critical set~{\cite[Definition~1.5]{CLS22}}]\label{def:critical-set}
Given a sweepout $\Phi$, define
\[
M(x) = \limsup_{r \to 0} \{\mathbf{M}(\Phi(y)): |y-x| < r\}.
\]
If $\Phi$ is an optimal sweepout, the \emph{critical domain of $\Phi$} is
\[
\mathfrak{m}(\Phi)=\{ x \in [0,1]: M(x) =W \}.
\]
We denote by $\mathfrak{m}_L(\Phi)$ and $\mathfrak{m}_R(\Phi)$ the sets of left and right critical points.
\end{definition}

\begin{definition}[Excessive intervals and points~{\cite[Definition~2.1]{CLS22}}]\label{def:excessive-interval}
Suppose $\{\Phi(x) = \partial\Omega(x)\}$ is a sweepout. Given a connected interval $I \subset [0,1]$
(which may be open, closed, or half-open), we say that $\{\Phi^I(x) = \partial\Omega^I(x)\}_{x\in \bar{I}}$
is an \emph{$I$-replacement family for $\Phi$} if
$\Omega^I(a) = \Omega(a)$, $\Omega^I(b) = \Omega(b)$ at the endpoints, and for all $x \in I$,
\[
\limsup_{I \ni y \to x} \mathbf{M}(\Phi^I(y)) < W,
\]
where $W := \sup_x \mathbf{M}(\Phi(x))$.

We say that a connected interval $I$ is an \emph{excessive interval for $\Phi$} if there exists an $I$-replacement family
for $\Phi$. A point $x$ is called \emph{left excessive for $\Phi$} (resp.\ \emph{right excessive}) if there exists an excessive interval $I$ with
$(x-\varepsilon, x] \subset I$ (resp.\ $[x, x+\varepsilon) \subset I$) for some $\varepsilon > 0$.
\end{definition}

A sweepout is called \emph{non-excessive} if no critical point is excessive on its critical side. The fundamental result of~\cite{CLS22} is that such sweepouts exist:

\begin{theorem}[Existence of non-excessive sweepouts~{\cite[Theorem~3.1]{CLS22}}]\label{thm:non-excessive-existence}
Let $(M^{n+1},g)$ be a closed Riemannian manifold with $2\leq n\leq 6$. There exists an (ONVP) sweepout $\Psi$ such that every $x\in \mathfrak{m}_L(\Psi)$ is not left excessive and every $x \in \mathfrak{m}_R(\Psi)$ is not right excessive.
\end{theorem}

We denote by $\Snm$ the collection of all non-excessive sweepouts. Any optimal sweepout produces a stationary varifold via the standard pull-tight argument:

\begin{proposition}[Stationary varifolds from optimal sweepouts~{\cite[Proposition~1.4]{CL03}}]\label{thm:CLS-stationary}
Let $(M^{n+1},g)$ be a closed Riemannian manifold with $n \geq 2$, and let $\Phi$ be an optimal sweepout with $\sup_x \mathbf{M}(\Phi(x)) = W$. Then there exists a stationary $n$-varifold $V$ in $M$ with $\mathbf{M}(V) \leq W$.
\end{proposition}

This follows from the pull-tight argument of Colding--de~Lellis~\cite{CL03}: varifold compactness produces a limit $V$ of slices whose masses approach $W$, and if every such limit were non-stationary, the tightening flow would produce a sweepout of strictly smaller width, contradicting optimality. The non-excessive condition is not used here; it enters in the analysis of the limit varifold below.

Let $\Phi$ be a non-excessive ONVP sweepout, $x_i \to x_0 \in \mathfrak{m}(\Phi)$ a min-max sequence, and $V$ the varifold limit of $|\partial^*\Omega(x_i)|$ given by Proposition~\ref{thm:CLS-stationary}. A key distinction arises depending on whether the mass of the limit slice equals the width:

\begin{definition}[Mass cancellation~{\cite[Section~6]{CLS22}}]\label{def:mass-cancellation}
Let $\Phi$ be a non-excessive (ONVP) sweepout with width $W$, and let $x_0 \in \mathfrak{m}(\Phi)$ be a critical point. We say that:
\begin{itemize}
    \item \emph{No mass cancellation} occurs at $x_0$ if $\mathbf{M}(\Phi(x_0)) = W$.
    \item \emph{Mass cancellation} occurs at $x_0$ if $\mathbf{M}(\Phi(x_0)) < W$.
\end{itemize}
\end{definition}

This dichotomy governs the integrality of the limit varifold; both cases are treated in Section~\ref{sec:integrality}.

Beyond stationarity and integrality, the non-excessive condition provides local geometric information about the limit varifold: critical slices are \emph{one-sided homotopic minimizers} in small balls, meaning they cannot be deformed to competitors of strictly smaller mass on either the inner or the outer side. This property is the mechanism through which non-excessiveness controls the regularity of the limit.

\begin{definition}[Homotopic minimizers~{\cite[Definition~1.4]{CLS22}}]\label{def:homotopic-minimizer}
Let $\Sigma = \bdy\Omega$ be a Caccioppoli set in $M$, and let $U \subset M$ be an open set. We say
that $\Sigma$ is:
\begin{itemize}
    \item \emph{inner homotopic minimizing in $U$} if there is no Caccioppoli set $E \subset \Omega$
    with $E \triangle \Omega \subset U$ and $\per(E, U) < \per(\Omega, U)$ that can be connected to
    $\Omega$ by a continuous family of Caccioppoli sets in $U$;
    \item \emph{outer homotopic minimizing in $U$} if there is no Caccioppoli set $E \supset \Omega$
    with $E \triangle \Omega \subset U$ and $\per(E, U) < \per(\Omega, U)$ that can be connected to
    $\Omega$ by a continuous family of Caccioppoli sets in $U$;
    \item \emph{one-sided homotopic minimizing in $U$} if it is either inner or outer homotopic
    minimizing in $U$.
\end{itemize}
\end{definition}

\begin{definition}[Non-homotopic-minimizing points~{\cite[Definition~2.5]{CLS22}}]\label{def:hnm}
For a varifold $V$ arising from a min-max procedure, we define the set of \emph{non-homotopic-minimizing
points} by
\[
    \mathfrak{h}_{\mathrm{nm}}(V) := \left\{P \in \spt\|V\| : \begin{array}{l}
        \text{for all $r > 0$ small, $\spt\|V\| \cap B_r(P)$} \\
        \text{is not one-sided homotopic minimizing in $B_r(P)$}
    \end{array}\right\}.
\]
\end{definition}

The set $\mathfrak{h}_{\mathrm{nm}}(V)$ consists of points where the varifold fails to be one-sided minimizing in every neighborhood. A key consequence of~\cite[Proposition~3.1]{CLS22} is that for non-excessive sweepouts, this set is at most finite. In particular, the limit varifold is one-sided homotopic minimizing---hence stable---at all but finitely many points. This will be used in Section~\ref{sec:regularity} to verify the hypotheses of Wickramasekera's regularity theorem.
